% !TeX program = xelatex
\documentclass[11pt,a4paper]{article}

\usepackage{polyglossia}
\setmainlanguage[numerals=arabic]{hebrew}
\setotherlanguage{english}
\newfontfamily\hebrewfont[Script=Hebrew,Language=Hebrew]{David CLM}
\newfontfamily\hebrewfonttt[Script=Hebrew,Language=Hebrew]{David CLM}
\newfontfamily\englishfont{Arial}

\usepackage{listings}
\usepackage{xcolor}
\usepackage{amsmath}
\usepackage{amssymb}
\usepackage[unicode]{hyperref}

\XeTeXgenerateactualtext 1 % Improve copy-paste and potentially fix some visual mirroring issues
\newcommand{\RLM}{\char"200F}
\newcommand{\LRM}{\char"200E}
\setcounter{secnumdepth}{-1} % Disable section numbering

\definecolor{dkgreen}{rgb}{0,0.6,0}
\definecolor{mauve}{rgb}{0.58,0,0.82}

\lstdefinelanguage{SMV}{
  keywords={MODULE, VAR, IVAR, ASSIGN, DEFINE, init, next, case, esac, TRUE, FALSE, boolean, INVARSPEC, LTLSPEC, SPEC},
  sensitive=true,
  comment=[l]{--},
  morecomment=[s]{/*}{*/},
  morestring=[b]",
}

\lstset{
  language=SMV,
  aboveskip=3mm,
  belowskip=3mm,
  showstringspaces=false,
  columns=fullflexible,
  keepspaces=true,
  basicstyle=\small\ttfamily,
  numbers=none,
  keywordstyle=\color{blue},
  commentstyle=\color{dkgreen},
  stringstyle=\color{mauve},
  breaklines=true,
  breakatwhitespace=true,
  tabsize=2
}

\newenvironment{englishlisting}
  {\begin{english}\begin{flushleft}}
  {\end{flushleft}\end{english}}

\renewcommand\labelitemi{--}

\title{תרגיל בית: מידול ואימות ב-\textenglish{SMV}}
\author{קורס שיטות אימות פורמליות}
\date{\today}

\begin{document}

\maketitle

\section{הנחיות הגשה}
יש להגיש קובץ דחוס אחד, המכיל \textbf{בדיוק} את ארבעת הקבצים הבאים:
\begin{itemize}
  \item \texttt{q1.smv}
  \item \texttt{q2.smv}
  \item \texttt{q3.smv}
  \item \texttt{answers.txt}
\end{itemize}

\noindent
הקובץ \texttt{answers.txt} חייב להיכתב \textbf{בדיוק} לפי התבנית שבסוף מסמך זה.

\subsection*{כללים מחייבים}
\begin{enumerate}
  \item בכל אחד מן הקבצים \texttt{q1.smv}, \texttt{q2.smv}, \texttt{q3.smv} חייב להיות בדיוק מודול עליון אחד בשם \texttt{main}.
  \item כאשר שאלה מגדירה במפורש שמות של משתנים, טיפוסים, או תחומי ערכים, אין לשנות אותם.
  \item מותר להוסיף \texttt{DEFINE} כרצונכם. משתני מצב או קלט נוספים מותרים רק אם השאלה אינה מקבעת ממשק חובה.
  \item כל מעבר חייב להיות מוגדר בצורה מלאה; אין להשאיר משתנה ללא הגדרת \texttt{next}.
  \item בכל שאלה יש לבדוק בדיוק את התכונות שנדרשו, ולא תכונות אחרות.
  \item ריצה, עקבת נגד או עדות נגישות שתוגש ב-\texttt{answers.txt} חייבת לציין \textbf{את כל המשתנים} של המודל בכל מצב.
\end{enumerate}

\section{שאלה 1: מציאת תקלה בפרוטוקול שגוי למניעה הדדית}
להלן פסאודו־קוד של שני תהליכים זהים, \texttt{P0} ו-\texttt{P1}, המשתמשים בשני דגלים משותפים:

\begin{englishlisting}
\begin{lstlisting}
shared boolean want0 := false;
shared boolean want1 := false;

process Pi (i in {0,1}, j = 1-i):
loop forever:
  L0: wait until !wantj;
  L1: wanti := true;
  L2: critical section;
  L3: wanti := false;
end loop
\end{lstlisting}
\end{englishlisting}

\subsection*{הנחות מידול מחייבות}
\begin{itemize}
  \item המערכת פועלת בשזירה: בכל צעד מתקדם תהליך אחד בלבד.
  \item כל שורה בפסאודו־קוד \textenglish{(L0, L1, L2, L3)} מייצגת צעד אטומי אחד.
  \item בפרט, בדיקת התנאי ב-\textenglish{L0} והתוצאה שלה (להישאר ב-\textenglish{L0} או לעבור ל-\textenglish{L1}) מתבצעות באותו צעד אטומי.
  \item הקטע הקריטי (L2) חייב להיות מיוצג במודל כמצב בקרה נפרד.
  \item בתחילת הריצה שני הדגלים הם \texttt{false}, ושני התהליכים מתחילים ב-\textenglish{L0}.
\end{itemize}

\subsection*{ממשק חובה}
על המודל להגדיר (דרך \texttt{VAR} או \texttt{DEFINE}) שני משתנים בוליאניים:
\begin{itemize}
  \item \texttt{in\_crit0} -- אמת אם ורק אם תהליך \texttt{P0} נמצא בקטע הקריטי.
  \item \texttt{in\_crit1} -- אמת אם ורק אם תהליך \texttt{P1} נמצא בקטע הקריטי.
\end{itemize}

\subsection*{התכונה שיש לבדוק}
יש להוסיף למודל \textbf{בדיוק} את תכונת הבטיחות הבאה:
\begin{englishlisting}
\begin{lstlisting}
INVARSPEC !(in_crit0 & in_crit1)
\end{lstlisting}
\end{englishlisting}

\subsection*{מה מותר ומה אסור}
\begin{itemize}
  \item חובה להשתמש בשמות \texttt{in\_crit0} ו-\texttt{in\_crit1} עבור הממשק.
  \item מותר לבחור כל שמות משתנים אחרים וכל קידוד סביר למצבי הבקרה.
  \item מותר להוסיף משתנה עזר לצורך שזירה, למשל משתנה בורר תהליך.
  \item אסור לשנות את התנהגות הפסאודו־קוד.
  \item אסור לפצל שורה אחת לשני צעדים אטומיים.
\end{itemize}

\subsection*{מה יש לענות}
ב-\texttt{answers.txt} יש למלא:
\begin{itemize}
  \item האם התכונה נכונה או שגויה.\RLM
  \item אם היא שגויה: אורך ריצת הנגד הקצרה ביותר.\RLM
  \item אם היא שגויה: ריצת הנגד הקצרה ביותר, כשרשומים בכל מצב כל משתני המודל.\RLM
  \item משפט קצר המסביר מהו הבאג ברמת האלגוריתם.\RLM
\end{itemize}

\section{שאלה 2: מציאת מרוץ בשינוי משתנה משותף}
להלן פסאודו־קוד של שני תהליכים, \texttt{P0} ו-\texttt{P1}, המבצעים פעולה על משתנה משותף \texttt{x}:

\begin{englishlisting}
\begin{lstlisting}
shared integer x := 1;

process Pi (i in {0,1}):
loop forever:
  L0: idle;
  L1: ok := (x > 0);
  L2: if (ok) then x := x - 1;
end loop
\end{lstlisting}
\end{englishlisting}

\subsection*{הנחות מידול מחייבות}
\begin{itemize}
  \item המערכת פועלת בשזירה: בכל צעד מתקדם תהליך אחד בלבד\RLM
  \item כל שורה בפסאודו־קוד \textenglish{(L0, L1, L2)} מייצגת צעד אטומי אחד\RLM
  \item בדיקת התנאי ב-\textenglish{L1} ועדכון המשתנה המקומי \texttt{ok} מתבצעים באותו צעד אטומי\RLM
  \item בתחילת הריצה \texttt{x=1}, ושני התהליכים מתחילים ב-\textenglish{L0}\RLM
\end{itemize}

\subsection*{ממשק חובה}
אין חובה על מבנה פנימי מסוים למודל, אך לצורך הגדירה והבדיקה חובה לעמוד בתנאים הבאים:
\begin{itemize}
  \item על המודל להכיל משתנה בשם \textenglish{x} בטווח \textenglish{-1..2}\RLM
  \item יש להוסיף למודל \textbf{בדיוק} את תכונת הבטיחות: \textenglish{INVARSPEC x >= 0}\RLM
  \item ב-\texttt{answers.txt} יש לציין את ערכי כל המשתנים במודל שבניתם (כולל משתני בקרה ודגלים מקומיים)\RLM
\end{itemize}

\subsection*{מה יש לענות}
ב-\texttt{answers.txt} יש למלא:
\begin{itemize}
  \item מהו הערך הקטן ביותר \texttt{v} שעבורו קיים מצב נגיש עם \texttt{x = v}.\RLM
  \item האם \texttt{INVARSPEC x >= 0} נכונה או שגויה.\RLM
  \item אורך הריצה הקצרה ביותר שמגיעה לראשונה לערך \texttt{v}.\RLM
  \item הריצה עצמה, כאשר בכל מצב מופיעים כל משתני המודל שבניתם.\RLM
\end{itemize}

\section{שאלה 3: מציאת שגיאה במימוש כפל של שני מספרים בני שני ביטים}
יש לבנות מודל \textenglish{SMV} למעגל קומבינטורי שגוי, המחשב מכפלה של שני מספרים בני שני ביטים.
המעגל אמור לקבל שני קלטים \texttt{a}, \texttt{b} בטווח \texttt{0..3}, ולחשב פלט \texttt{out} בטווח \texttt{0..15}.

\subsection*{ממשק חובה}
הקובץ \texttt{q3.smv} חייב להתחיל במבנה הבא:
\begin{englishlisting}
\begin{lstlisting}
MODULE main
VAR
  a1 : boolean;
  a0 : boolean;
  b1 : boolean;
  b0 : boolean;
\end{lstlisting}
\end{englishlisting}

\subsection*{סמנטיקה של הקלטים}
\begin{itemize}
  \item הקלטים צריכים להיות לא-דטרמיניסטיים במצב ההתחלתי\RLM
  \item לאחר האתחול, עליהם להישאר קבועים לנצח\RLM
\end{itemize}

\subsection*{הגדרות חובה}
יש להגדיר באמצעות \texttt{DEFINE} את הגדלים הבאים:
\begin{itemize}
  \item \texttt{a = 2*a1 + a0}
  \item \texttt{b = 2*b1 + b0}
  \item \texttt{p0 = a0 \& b0}
  \item \texttt{p1 = a1 \& b0}
  \item \texttt{p2 = a0 \& b1}
  \item \texttt{p3 = a1 \& b1}
  \item \texttt{r0 = p0}
  \item \texttt{r1 = p1 xor p2}
  \item \texttt{c1 = p1 \& p2}
  \item \texttt{r2 = p3 \& c1}
  \item \texttt{r3 = p3 \& c1}
  \item \texttt{out = 8*r3 + 4*r2 + 2*r1 + r0}
\end{itemize}

\subsection*{התכונה שיש לבדוק}
יש להוסיף למודל את התכונה:
\begin{englishlisting}
\begin{lstlisting}
INVARSPEC out = a * b
\end{lstlisting}
\end{englishlisting}

\subsection*{כלל למינימליות}
אם התכונה שגויה, הזוג הכושל המינימלי ייקבע לפי סדר לקסיקוגרפי:
\begin{itemize}
  \item תחילה ממזערים את \texttt{a}\RLM
  \item מבין כל הזוגות עם אותו \texttt{a}, ממזערים את \texttt{b}\RLM
\end{itemize}

\subsection*{מה יש לענות}
ב-\texttt{answers.txt} יש למלא:
\begin{itemize}
  \item האם התכונה נכונה או שגויה.\RLM
  \item אם היא שגויה: הזוג הכושל המינימלי \textenglish{(a,b)}.\RLM
  \item ארבעת ביטי הקלט המתאימים לזוג זה.\RLM
  \item הערך הצפוי \texttt{a*b}.\RLM
  \item הערך בפועל \texttt{out}.\RLM
  \item איזו שורה שגויה: \texttt{r1}, \texttt{r2}, או \texttt{r3}.\RLM
  \item תיקון לשורה השגויה בלבד.\RLM
\end{itemize}

\section{תבנית חובה לקובץ \texttt{answers.txt}}
יש להגיש את הקובץ \texttt{answers.txt} בדיוק בתבנית הבאה:

\begin{englishlisting}
\begin{lstlisting}
[Q1]
Verdict = TRUE/FALSE
ShortestCounterexampleLength = <integer or N/A>
State0 = (...)
State1 = (...)
...
BugExplanation = <one short sentence>

[Q2]
MinReachableX = <integer>
Invariant = TRUE/FALSE
ShortestWitnessLength = <integer>
State0 = (...)
State1 = (...)
...

[Q3]
Verdict = TRUE/FALSE
SmallestFailingPair = (a=<integer>, b=<integer>)
Bits = (a1=<0/1>, a0=<0/1>, b1=<0/1>, b0=<0/1>)
Expected = <integer>
Actual = <integer>
WrongLine = r1/r2/r3
Fix = <equation or N/A>
\end{lstlisting}
\end{englishlisting}

\noindent
הערות:
\begin{itemize}
  \item אם בשאלה 1 התכונה נכונה, כתבו \texttt{ShortestCounterexampleLength = N/A} ואל תוסיפו מצבים.
  \item בשאלה 1, בכל מצב בריצת הנגד יש לציין את כל משתני המודל שבחרתם.
  \item אם בשאלה 3 התכונה נכונה, כתבו \texttt{Fix = N/A}.
  \item בכל ריצה או עדות נגישות יש לציין את כל משתני המודל בכל מצב.
\end{itemize}

\end{document}
